\subsection{What can actually be rebuilt with fewer even lift parameters}

\begin{proposition}[DP44: smaller rank records and the exact remaining comparison]\label{DP44}
Let $R$ be a recovered \path{Ranks} record with lift count $M$, fixed frame, and parameters $k,t$. For every even $M'$ with $2\le M'\le M$:
\begin{enumerate}[label=(\alph*),leftmargin=*]
\item Replacing $M$ by $M'$ gives a valid rank record. Every \path{HCParam R} restricts by keeping $u,v$ and the first $M'$ entries of $b,e$.
\item Every \path{Host R} transports with the same fields and a weaker size proof. Given the same prime $T$ and norm element $z\bar z=T$, the source's actual generic constructor supplies the new Hermitian \path{Spaces}.
\item At fixed $k,t$, retained positive spaces $X_j$ keep their standard forms and ranks $k+j(t+1)$, but the auxiliary signatures and terminal period change. This reconstruction is not literal truncation with the old endpoint.
\item For the source's gain and main-term definitions, if the other data are held fixed, put
\[
 g_j=\sum_w\operatorname{sz}_w(\operatorname{gainD}_{j,w}),\quad
 e_j=\sum_w\operatorname{sz}_w(E_{u,j,w}),\quad
 T_\tau=\sum_w\operatorname{sz}_w(\tau_w),
\]
\[
 H=\sum_t(\operatorname{phF}(A_t)-A_t\log\Omega),\qquad
 K_u=\sum_t A_t\log T_u.
\]
Then
\[
 \operatorname{Ga}_m=\sum_{j=1}^m g_j-T_\tau,\qquad
 \operatorname{Ee}_m=\sum_{j=1}^m e_j,\qquad
 \operatorname{MAIN}_m=mH-K_u.
\]
In particular, with $Q_m=\operatorname{Ga}_m-\operatorname{Ee}_m-\operatorname{MAIN}_m$,
\[
 Q_{M'}-Q_M=-\sum_{j=M'+1}^{M}(g_j-e_j-H).
\]
If an allowance $B_m$ is also subtracted, the right side acquires $B_M-B_{M'}$. For genuinely rebuilt families, every changed input must instead be included in the comparison.
\end{enumerate}
The proposition neither constructs a nonzero shorter theta family with its required estimates nor asserts minimality of the old lift count.
\end{proposition}

\begin{proof}
The full \path{Ranks} fields are in \path{Formalism/LocalRingCardinality/Params.lean}, lines 6--22, and its \path{HCParam} fields at 50--69. The inequalities
\[
 M'+3\le k,\quad k+M'-1+5\le t,\quad 3M'<t,\quad
 M'\le2\operatorname{ord}\le k
\]
follow from their old versions. The parity assumptions on $k,t$ and the frame are unchanged. The remaining nonlinear field follows from $M'\le t+1$:
$(M')^2\le M'(t+1)\le k+M'(t+1)$.
This proves every rank requirement.
The interlacing conditions on $u,v$ and the first $b$ are unchanged; every retained inequality
$e_j-b_j\ge(t+3)/2$ and $b_{j+1}-e_j\ge(t+3)/2$ is an old one. Integral and half-integral entries remain so. The same restriction preserves the word-family conditions in \path{SpanPowerMembership/NZRowAsm.lean}, lines 15--29.

The full \path{Host} structure in \path{DataHost.lean}, lines 39--109, depends on the changed $m$ only through its size condition
$4t(k+m-1)+1<q,d$ at line 54. That condition weakens when $M$ is replaced by $M'$. All other fields use the same frame, $k,t$, or independent data.
The generic \path{hostSpacesOfNorm} in
\path{SpanPowerMembership/RankRecordRowDimensionZero.lean}, lines 137--142, calls \path{spacesOfRanks}, using the same host's \path{cmData}, integral involution and \path{exists_adelic_neg_one}, together with the assumed norm equation.
The actual constructor is \path{SpanPowerMembership/BaseFormAnisotropicOrthogonalSingle.lean}, hypotheses at 200--205 and \path{spacesOfRanks} at 325--386. It constructs $X_j$ as the positive standard space of rank $k+j(t+1)$, $Y_j$ of rank $k+M'-1+jt$ with negative index $\floor{(M'-j)/2}$, and $W_j$ of rank $k+M'-1$ with signature $(k+j-1,M'-j)$; its distinguished vectors, complement and adelic isometry fields are filled in the cited constructor. Thus the new spaces exist under the stated inputs.

Their retained $X_j$ are unchanged, but
\[
 N_M-N_{M'}=(M-M')(t+1)>0\quad(M'<M).
\]
The old $W_{M'}$ has signature $(k+M'-1,M-M')$; the new terminal $W'_{M'}$ has signature $(k+M'-1,0)$. This proves the endpoint qualification.
Even equality of terminal positive dimensions alone does not force the lift count if $t$ changes: given even $2\le M'<M$, even $k\ge M+3$ and a frame with $M\le2\operatorname{ord}\le k$, choose a sufficiently large integer $c$ and put
$t+1=2M'c$, $t'+1=2Mc$.
Both rank records then satisfy all their inequalities and
$k+M(t+1)=k+M'(t'+1)$.
This numerical construction transports no host bounds, gaps, signatures or period estimates.

Finally \path{DataParams}\textsubscript{2}\path{.lean} defines \path{sz} at 203--206 as a scalar multiple, the lift set at 208, \path{gainD} at 212--213, \path{Fo/Ga/Ee} at 218--222, and \path{MAIN} at 232--234. Distributing the finite sums gives the three displayed identities. Subtracting them for $M,M'$ gives the identity for $Q$ and the optional allowance. No sign assumption on $H$ was used. This proves all parts.
\end{proof}

The next actual input is now precise. The generic nonzero-row constructor
\path{SpanPowerMembership/NZRowE.lean}, \path{hostLevelRecord} at line 194, uses the generic spaces at 198, but its hypotheses at 165--189 include class lattices, a level window, norm congruences, signature weights, local gap and $G$ data, and at 185 a period
\[
 \operatorname{per}_u:\operatorname{FK}(\operatorname{cls}_u(M))
        \longrightarrow\operatorname{En}(d)\longrightarrow\mathbf C.
\]
At $M'$ that period has a different domain. Composing the old period with discarded kernels still uses those kernels; no cheaper geometric operation has thereby been exhibited. The similarly named \path{rankRecord} in the constructor source has empty class sets at lines 25--44 and zero period at 77--83, so it is not a nonzero shorter-chain witness.

The quantitative inputs remain explicit as well.
\path{NZRowAsm.lean}, \path{ExtraNZIn}, lines 61--86, requires the level-error and lattice comparisons, actual local-to-global kernels, entry and entropy ties, the $\operatorname{Kap}$ data, orbit-filling data, normalized dimensions and gauges. Its \path{extraNZ}, lines 91--123, assembles these inputs; it does not eliminate them. \path{GlobalTheoremInputs.lean}, lines 158--160, then requires \path{LinkG}, \path{LinkS}, \path{LinkIRR}, \path{LinkDIV} and \path{LinkMu} for the new family.

In particular \path{LinkG_at}, lines 122--136, is the one-way lower certificate
\[
 \sum_{j=1}^m\sum_{t\in\mathrm{TT}_w}
       (R_{2,j,t,p}-\operatorname{loG})-m\operatorname{flG}
 \le\sum_{j=1}^m\operatorname{gain}_{j,w}-\tau_w.
\]
The odd-level ramp $R_2$ is \path{vonMangoldtClassSum},
\path{DataParams}\textsubscript{2}\path{.lean}, lines 240--242, with cutoff
$\operatorname{Nat.log}_p(a_T^{\mathrm{toNat}})+2$ and weights at least one. This lower certificate cannot be reversed into an upper bound on actual gain.
The upper comparison \path{LinkDIV_at}, lines 77--89, contains
$\operatorname{MAIN}$, $m(B_0+N_c)$ and further costs; the local
\path{LocalQuadraticParts.lean}, \path{Conclusion4}, lines 677--680, uses
$\operatorname{Ga}+\operatorname{Fo}-\operatorname{Ee}$.
Other quantities also depend on $m$: \path{Gam_eq} at
\path{DataParams}\textsubscript{2}\path{.lean}, line 176, is
$C_\Gamma m^5fL$, and \path{TuGe} at 252 requires $b^{m/2}\le T$.
\path{flG}, lines 100--102, and \path{loG}, lines 109--120 of
\path{GlobalTheoremInputs.lean}, contain terminal ranks and further orbit/scale costs.

Thus rank/HC/host compatibility and generic Hermitian spaces are not the remaining obstruction for a smaller even parameter in the existing host. The next gate is a new endpoint period with its ties, followed by the actual selected local inputs and net cost comparison. A mixed CM host additionally has to replace DP25's induced-type input. DP40's equal-direction sum and DP42's primitive-line rates settle their stated local families, but neither identifies the full selected determinants with those rates or reverses the source's gain certificate.

